% Options for packages loaded elsewhere
\PassOptionsToPackage{unicode}{hyperref}
\PassOptionsToPackage{hyphens}{url}
\PassOptionsToPackage{dvipsnames,svgnames,x11names}{xcolor}
\documentclass[
  10pt,
  fontset=fandol]{ctexart}
\usepackage{xcolor}
\usepackage[margin=16mm]{geometry}
\usepackage{amsmath,amssymb}
\setcounter{secnumdepth}{-\maxdimen} % remove section numbering
\usepackage{iftex}
\ifPDFTeX
  \usepackage[T1]{fontenc}
  \usepackage[utf8]{inputenc}
  \usepackage{textcomp} % provide euro and other symbols
\else % if luatex or xetex
  \usepackage{unicode-math} % this also loads fontspec
  \defaultfontfeatures{Scale=MatchLowercase}
  \defaultfontfeatures[\rmfamily]{Ligatures=TeX,Scale=1}
\fi
\usepackage{lmodern}
\ifPDFTeX\else
  % xetex/luatex font selection
\fi
% Use upquote if available, for straight quotes in verbatim environments
\IfFileExists{upquote.sty}{\usepackage{upquote}}{}
\IfFileExists{microtype.sty}{% use microtype if available
  \usepackage[]{microtype}
  \UseMicrotypeSet[protrusion]{basicmath} % disable protrusion for tt fonts
}{}
\makeatletter
\@ifundefined{KOMAClassName}{% if non-KOMA class
  \IfFileExists{parskip.sty}{%
    \usepackage{parskip}
  }{% else
    \setlength{\parindent}{0pt}
    \setlength{\parskip}{6pt plus 2pt minus 1pt}}
}{% if KOMA class
  \KOMAoptions{parskip=half}}
\makeatother
\setlength{\emergencystretch}{3em} % prevent overfull lines
\providecommand{\tightlist}{%
  \setlength{\itemsep}{0pt}\setlength{\parskip}{0pt}}
\usepackage{amsmath,amssymb,mathtools,needspace}
\xeCJKDeclareCharClass{CJK}{"2032,"0400->"04FF,"2160->"216B,"2460->"2473,"25A0,"25C7,"25CB,"25CF}
\IfFontExistsTF{Arial Unicode MS}{\xeCJKsetup{AutoFallBack=true}\setCJKfallbackfamilyfont{\CJKrmdefault}{Arial Unicode MS}}{}
\setcounter{secnumdepth}{0}
\setlength{\emergencystretch}{2em}
\allowdisplaybreaks
\usepackage{bookmark}
\IfFileExists{xurl.sty}{\usepackage{xurl}}{} % add URL line breaks if available
\urlstyle{same}
\hypersetup{
  pdftitle={基于数值积分的构造方法},
  colorlinks=true,
  linkcolor={Maroon},
  filecolor={Maroon},
  citecolor={Blue},
  urlcolor={Blue},
  pdfcreator={LaTeX via pandoc}}

\title{基于数值积分的构造方法}
\author{}
\date{}

\begin{document}
\maketitle

\subsubsection{5.5.1 基于数值积分的构造方法}\label{ux57faux4e8eux6570ux503cux79efux5206ux7684ux6784ux9020ux65b9ux6cd5}

将方程 \(y'=f(x,y)\) 的两端从 \(x_n\) 到 \(x_{n+1}\) 求积分，得

\[
y(x_{n+1})=y(x_n)+\int_{x_n}^{x_{n+1}}f(x,y(x))\,\mathrm dx.
\tag{5.5.1}
\]

为了通过这个积分关系式获得 \(y(x_{n+1})\) 的近似值，只要近似地算出其中的积分项 \(\int_{x_n}^{x_{n+1}}f(x,y(x))\,\mathrm dx\) 即可。选用不同的数值方法计算这个积分项，就会导出不同的计算格式。

例如，设用矩形方法计算积分项，即

\[
\int_{x_n}^{x_{n+1}}f(x,y(x))\,\mathrm dx\approx hf(x_n,y(x_n)),
\]

代入式（5.5.1），有

\[
y(x_{n+1})\approx y(x_n)+hf(x_n,y(x_n)),
\]

据此离散化即可导出 Euler 格式。

\href{../../source-images/pdf-138.jpeg}{原页}

为了改善精度，可以改用梯形方法计算积分项，即

\[
\int_{x_n}^{x_{n+1}}f(x,y(x))\,\mathrm dx\approx\frac h2[f(x_n,y(x_n))+f(x_{n+1},y(x_{n+1}))],
\]

再代入式（5.5.1），有

\[
y(x_{n+1})\approx y(x_n)+\frac h2[f(x_n,y(x_n))+f(x_{n+1},y(x_{n+1}))],
\]

据此离散化即可导出格式（5.2.8），\textbf{梯形法}因此而得名。

我们知道，基于插值原理可以建立一系列的数值积分方法，运用这些方法可以导出求解微分方程的一系列的计算格式。一般地，设已构造出 \(f(x,y(x))\) 的插值多项式 \(P_r(x)\)，那么，计算 \(\int_{x_n}^{x_{n+1}}P_r(x)\,\mathrm dx\) 作为 \(\int_{x_n}^{x_{n+1}}f(x,y(x))\,\mathrm dx\) 的近似值，即可将式（5.5.1）离散化得到下列计算公式：

\[
y_{n+1}=y_n+\int_{x_n}^{x_{n+1}}P_r(x)\,\mathrm dx.
\tag{5.5.2}
\]

\begin{center}\rule{0.5\linewidth}{0.5pt}\end{center}

\subsection{相关原页校注（agent 补充）}\label{ux76f8ux5173ux539fux9875ux6821ux6ce8agent-ux8865ux5145}

以下校注来自本节覆盖的原页，可能也涉及同页相邻小节；原文照录与校正说明分开保留。

\subsubsection{原书 PDF 页 138 的校注}\label{ux539fux4e66-pdf-ux9875-138-ux7684ux6821ux6ce8}

\href{../pages/pdf-138.md}{完整原页转录} · \href{../../source-images/pdf-138.jpeg}{原页图像}

校注记录：ch05b-E001。

\begin{quote}
校注（agent 补充）：原图中上述未编号``差分展开式''的求和上限印作 \(r\)，求和项末尾印作 \(f_{n-j}\)，这里均保留原文。末尾下标疑为原书排印错误，通常应为 \(f_{n-i}\)；例如 \(j=r=1\) 时，原印刷式右边为零，而左边为 \(f_n-f_{n-1}\)。按通常写法求和上限为 \(j\)；若约定 \(\binom ji=0\)（\(i>j\)），则上限写 \(r\) 也可。此疑误已对照原页局部图，并非 OCR 替换。
\end{quote}

\subsection{本节覆盖原页的校注记录}\label{ux672cux8282ux8986ux76d6ux539fux9875ux7684ux6821ux6ce8ux8bb0ux5f55}

下列记录按原页关联，含同页相邻小节的校注；计算时同时读取上文校注和完整原页。

\begin{itemize}
\tightlist
\item
  \texttt{ch05b-E001}：原书 PDF 页 138
\end{itemize}

\end{document}
